%----------
%	CONCLUSION OF THESIS
%---------------------------------------------------------------------------------

\addchaptertocentry{Дүгнэлт}
\pagecolor{white}

\section*{\centering \huge{Дүгнэлт}}

Энэхүү бакалаврын төгсөлтийн ажлын хүрээнд Өмнөговь аймгийн Даланзадгад хотын \textbf{Өмнөговь аймгийн Спорт цогцолбор}-т зориулсан \textbf{үүлэн орчинд суурилсан цаг захиалгын мобайл аппликейшн}-ийг Flutter, Node.js/Express, Firebase, Docker, AWS болон AI технологиудыг хосолсон орчин үеийн бүрэн стекэд хөгжүүлж, дараахи зорилтуудыг хэрэгжүүлсэн:

\begin{enumerate}
    \item \textbf{Асуудлын судалгаа ба шаардлагын шинжилгээ.} Даланзадгад хотын иргэдийн тулгамдаж буй таван үндсэн асуудлыг тогтоож — утсаар захиалах, биечлэн очих, цагийн мэдээлэл ил тод бус, цуцлах боломжгүй, сануулга байхгүй гэсэн --- тус бүрд нь цахим шийдлийг тодорхойлсон. Функциональ болон функциональ бус шаардлагуудыг нарийвчлан баталгаажуулсан.

    \item \textbf{Flutter мобайл аппликейшн.} Material Design 3 стандартад нийцсэн, дулаан говийн цайвар дизайн (бежийн дэвсгэр \code{\#EDE7DA}, наранжин \code{\#D4700F}, Montserrat фонт) бүхий мобайл аппликейшнийг 4 tab-тай доод навигацитай хэрэгжүүлсэн. StatefulWidget + setState суурилсан төлвийн удирдлага нэмэлт dependency байхгүйгээр тогтвортой ажилласан. Апп нь Android болон iOS платформд нэгэн зэрэг ажилладаг.

    \item \textbf{Firebase үйлчилгээ ба email нэвтрэлт.} Firebase Authentication-ийн email/нууц үгийн нэвтрэлтийг Node.js бэкэнд дамжуулан Nodemailer (Gmail SMTP)-ийн 6 оронтой баталгаажуулах кодтой хосолсон. Кодын хугацаа 10 минут бөгөөд амжилтгүй тохиолдолд Firebase rollback хэрэгжинэ. Cloud Firestore-д \code{bookings}, \code{users}, \code{fixed\_bookings}, \code{email\_verifications} гэсэн дөрвөн collection тодорхойлж, DBK-xxxxx захиалгын кодын систем хэрэгжүүлсэн.

    \item \textbf{Node.js/Express REST API бэкэнд.} 8 endpoint бүхий RESTful API хэрэгжүүлж, email баталгаажуулалт, захиалгын менежмент, цагийн хуваарь, AI чат зэрэг функцийг хангасан. \code{lib/config.dart} файлаар хөгжүүлэлтийн болон продакшн орчны URL-ийг нэг газраас удирддаг.

    \item \textbf{Үүлэн байршуулалт: Docker + AWS EC2.} Бэкэнд сервер нь Docker контейнержуулалтаар \textbf{AWS EC2 (ap-southeast-1, Сингапур)} дээр бодитоор ажиллаж байна. AWS ECR-д Docker image хадгалж, EC2 дээр ажиллуулдаг энэ cloud-native архитектур нь системийн тасралтгүй байдал, хэмжээний уян хатан байдлыг хангадаг.

    \item \textbf{GitHub Actions CI/CD автоматжуулалт.} \code{backend/} хавтаст өөрчлөлт push хийхэд автоматаар Docker image build → AWS ECR push → EC2 SSH deploy гэсэн дарааллаар 3--4 минутын дотор байршуулалт хийгддэг. 5 дараалсан push-д туршиж, тасралтгүй ажилласныг баталгаажуулсан.

    \item \textbf{Давхар горимын AI чатбот.} Anthropic Claude Haiku 4.5-г Firebase Cloud Functions дамжуулан (үндсэн горим) болон Groq Llama 3.1-8b-г Node.js бэкэнд дамжуулан (нөөц горим) ажиллуулах нөөцлөлтийн архитектур нэвтрүүлсэн. Энэ нь AI үйлчилгээний найдвартай байдлыг нэмэгдүүлнэ.

    \item \textbf{Туршилтын үр дүн.} Захиалгын бүрэн мөчлөгийн функциональ туршилтын 9 алхам (бүртгэлээс эхлэн захиалга цуцлах, давхардал шалгалт хүртэл) бүгд амжилттай давсан. Гүйцэтгэлийн хэмжилтийн 5 үзүүлэлт (апп ачаалах, Firestore хүсэлт, EC2 API, захиалга үүсгэх, AI чатбот) тус бүр шаардлагын хязгаарт нийцсэн. Апп 2.3 секундын дотор ачаалж, Firestore хүсэлт 0.6 с, EC2 API 0.8 с хариулдаг болохыг баталгаажуулсан.
\end{enumerate}

\textbf{Туршилтын нэгтгэсэн дүн:}

\begin{table}[h]
\centering
\caption{Зорилтуудын биелэлт}
\label{tab:goals}
\begin{tabularx}{\textwidth}{|p{0.4cm}|X|p{2.5cm}|}
\hline
\textbf{\#} & \textbf{Зорилт} & \textbf{Биелэлт} \\
\hline
1 & Асуудлын судалгаа, шаардлагын шинжилгээ & \textcolor{green!60!black}{\textbf{Бүрэн}} \\
\hline
2 & UML диаграммууд (Use Case, Activity, Sequence) & \textcolor{green!60!black}{\textbf{Бүрэн}} \\
\hline
3 & Flutter мобайл аппликейшн (Android/iOS) & \textcolor{green!60!black}{\textbf{Бүрэн}} \\
\hline
4 & Firebase Auth + Firestore + Cloud Functions & \textcolor{green!60!black}{\textbf{Бүрэн}} \\
\hline
5 & Node.js/Express REST API & \textcolor{green!60!black}{\textbf{Бүрэн}} \\
\hline
6 & Docker + AWS EC2 үүлэн байршуулалт & \textcolor{green!60!black}{\textbf{Бүрэн}} \\
\hline
7 & GitHub Actions CI/CD & \textcolor{green!60!black}{\textbf{Бүрэн}} \\
\hline
8 & AI чатбот (Claude Haiku + Groq Llama) & \textcolor{green!60!black}{\textbf{Бүрэн}} \\
\hline
9 & Функциональ болон гүйцэтгэлийн туршилт & \textcolor{green!60!black}{\textbf{Бүрэн}} \\
\hline
\end{tabularx}
\end{table}

\textbf{Цаашдын хөгжлийн чиглэл:}

\begin{itemize}
    \item \textbf{Онлайн төлбөр} — QPay, SocialPay зэрэг Монголын төлбөрийн системүүдтэй нэгтгэх;
    \item \textbf{Администраторын дэлгэц} — Спорт заалны ажилтнуудад зориулсан хуваарь удирдах вэб дэлгэц хөгжүүлэх;
    \item \textbf{HTTPS/SSL} — EC2 болон Flutter апп хоорондын холболтод SSL нэмж аюулгүй байдлыг сайжруулах;
    \item \textbf{Firebase Storage} — Спорт заалнуудын зургийг Firestore-тэй уялдуулан хадгалах;
    \item \textbf{GPS байршил} — Хэрэглэгчид хамгийн ойр заалыг харуулах функц;
    \item \textbf{Статистик самбар} — Захиалгын дүн шинжилгээ, ачаалалтай цагийн мэдээлэл;
    \item \textbf{Өргөтгөл} — Даланзадгад хотоос гадна Өмнөговь аймгийн бусад сумдын спорт байгууламжид өргөтгөх.
\end{itemize}

Энэхүү ажил нь Монгол улсын орон нутгийн хотуудад ухаалаг цахим үйлчилгээ нэвтрүүлэх боломжийг тодорхой харуулж, Flutter болон Firebase-г Node.js/Express REST API, Docker/AWS үүлэн дэд бүтэц, GitHub Actions CI/CD автоматжуулалт, AI чатбот үйлчилгээтэй хосолсон бүрэн стекийн туршлагыг цуглуулсан чухал бүтээл болсон гэж дүгнэж байна. Хөгжүүлсэн систем нь бодитоор үүлэн орчинд (AWS EC2) ажиллаж байгаа нь ирээдүйн орон нутгийн цахим үйлчилгээний системүүдэд загвар болж чадна.
